% =========================================================
% Section: Derivative Geometry --- Extrema
% =========================================================

\begin{tcolorbox}[colback=gray!8, colframe=gray!40, title=\textbf{Derivative Geometry --- toolkit}]
\textbf{Concept family:} Relative extrema, the Interior Extremum Theorem, and the necessary condition for extrema.

\medskip
\textbf{Atomic items in this section:}
\begin{itemize}
  \item Definition: Relative Minimum
  \item Definition: Relative Maximum
  \item Definition: Relative Extremum
  \item Theorem: Interior Extremum Theorem
  \item Theorem: Necessary Condition for an Interior Extremum
  \item Corollary: Necessary Condition for an Extremum
\end{itemize}
\end{tcolorbox}

% ---------------------------------------------------------

\begin{tcolorbox}[colback=corbox, colframe=corborder]
\begin{definition}[Relative Minimum]
\label{def:relative-minimum}
Let $A \subseteq \mathbb{R}$, let $f : A \to \mathbb{R}$, and let $c \in A$. We say that $f$ has a \emph{relative minimum} at $c$ if there exists $\delta > 0$ such that
\[
f(c) \leq f(x) \qquad \text{for all } x \in A \cap (c-\delta,c+\delta).
\]
If, in addition,
\[
f(c) < f(x) \qquad \text{for all } x \in \bigl(A \cap (c-\delta,c+\delta)\bigr)\setminus\{c\},
\]
then $f$ has a \emph{strict relative minimum} at $c$.
\end{definition}
\end{tcolorbox}

\begin{remark*}[Standard quantified statement]
\[
\exists \delta > 0\;\forall x \in A\;
\bigl(|x - c| < \delta \implies f(c) \leq f(x)\bigr).
\]
\end{remark*}

\begin{remark*}[Definition predicate reading]
\[
\operatorname{LocalMinimum}(f, c)
\;\Longleftrightarrow\;
\exists \delta > 0\;\forall x \in A\;
\bigl(\operatorname{Within}(x, c, \delta) \implies f(c) \leq f(x)\bigr).
\]
\end{remark*}

\begin{remark*}[Interpretation]
$f$ has a relative minimum at $c$ if $f(c)$ is the smallest value of $f$ in some open neighbourhood of $c$ within the domain. The word \emph{relative} (or \emph{local}) signals that we are comparing $f(c)$ only to nearby values, not to all values of $f$ globally. The strict variant excludes the degenerate case where nearby points share the minimum value. The symmetric notion --- relative maximum --- reverses the inequality. Together they constitute a relative extremum. These are purely neighbourhood conditions; no derivative appears in the definitions.
\end{remark*}

\NoLocalDependencies
% ---------------------------------------------------------

\begin{tcolorbox}[colback=corbox, colframe=corborder]
\begin{definition}[Relative Maximum]
\label{def:relative-maximum}
Let $A \subseteq \mathbb{R}$, let $f : A \to \mathbb{R}$, and let $c \in A$. We say that $f$ has a \emph{relative maximum} at $c$ if there exists $\delta > 0$ such that
\[
f(x) \leq f(c) \qquad \text{for all } x \in A \cap (c-\delta,c+\delta).
\]
If, in addition,
\[
f(x) < f(c) \qquad \text{for all } x \in \bigl(A \cap (c-\delta,c+\delta)\bigr)\setminus\{c\},
\]
then $f$ has a \emph{strict relative maximum} at $c$.
\end{definition}
\end{tcolorbox}

\begin{remark*}[Standard quantified statement]
\[
\exists \delta > 0\;\forall x \in A\;
\bigl(|x - c| < \delta \implies f(x) \leq f(c)\bigr).
\]
\end{remark*}

\begin{remark*}[Definition predicate reading]
\[
\operatorname{LocalMaximum}(f, c)
\;\Longleftrightarrow\;
\exists \delta > 0\;\forall x \in A\;
\bigl(\operatorname{Within}(x, c, \delta) \implies f(x) \leq f(c)\bigr).
\]
\end{remark*}

\begin{remark*}[Dependencies]
\begin{itemize}
  \item \hyperref[def:relative-minimum]{Definition: Relative Minimum}
\end{itemize}
\end{remark*}

% ---------------------------------------------------------

\begin{tcolorbox}[colback=corbox, colframe=corborder]
\begin{definition}[Relative Extremum]
\label{def:relative-extremum}
Let $A \subseteq \mathbb{R}$, let $f : A \to \mathbb{R}$, and let $c \in A$. We say that $f$ has a \emph{relative extremum} at $c$ if $f$ has either a relative minimum at $c$ or a relative maximum at $c$.
\end{definition}
\end{tcolorbox}

\begin{remark*}[Standard quantified statement]
\[
\operatorname{LocalExtremum}(f, c)
\;\Longleftrightarrow\;
\operatorname{LocalMinimum}(f, c) \;\lor\; \operatorname{LocalMaximum}(f, c).
\]
\end{remark*}

\begin{remark*}[Definition predicate reading]
\[
\operatorname{LocalExtremum}(f, c)
\;\Longleftrightarrow\;
\operatorname{LocalMinimum}(f, c) \;\lor\; \operatorname{LocalMaximum}(f, c).
\]
\end{remark*}

\begin{remark*}[Dependencies]
\begin{itemize}
  \item \hyperref[def:relative-minimum]{Definition: Relative Minimum}
  \item \hyperref[def:relative-maximum]{Definition: Relative Maximum}
\end{itemize}
\end{remark*}

% ---------------------------------------------------------

\begin{theorem}[Interior Extremum Theorem]
\label{thm:interior-extremum-theorem}
Let $I \subseteq \mathbb{R}$ be an open interval, let $f : I \to \mathbb{R}$, and let $c \in I$.
If $f$ has a relative extremum at $c$ and $f$ is differentiable at $c$, then
\[
f'(c) = 0.
\]
\hyperref[prf:interior-extremum-theorem]{Go to proof.}
\end{theorem}

\begin{remark*}[Standard quantified statement]
\[
\operatorname{LocalExtremum}(f, c) \;\land\; \operatorname{DifferentiableAt}(f, c)
\;\Rightarrow\; \operatorname{DerivativeAt}(f, c, 0).
\]
\end{remark*}

\begin{remark*}[Contrapositive quantified statement]
\[
\operatorname{DerivativeAt}(f, c, L) \;\land\; L \neq 0
\;\Rightarrow\; \lnot\operatorname{LocalExtremum}(f, c).
\]
If $f'(c)$ exists and is nonzero, then $c$ is not a relative extremum.
\end{remark*}

\begin{remark*}[Contrapositive predicate reading]
\[
\operatorname{DerivativeAt}(f, c, L) \;\land\; L \neq 0
\;\Rightarrow\;
\lnot\operatorname{LocalMinimum}(f, c) \;\land\; \lnot\operatorname{LocalMaximum}(f, c).
\]
\end{remark*}

\begin{remark*}[Interpretation]
The theorem gives a necessary condition for interior extrema under differentiability: the derivative must vanish. The geometric picture is immediate --- at a local max or min, the tangent line must be horizontal, since the function is neither consistently increasing nor consistently decreasing. The proof is a sign argument: at a relative maximum, the difference quotient is non-positive for $x > c$ and non-negative for $x < c$; since the derivative is a single limit, it is squeezed to zero. The converse fails: $f'(c) = 0$ does not imply an extremum ($f(x) = x^3$ at $c = 0$). Critical points --- where $f'(c) = 0$ or $f'(c)$ does not exist --- are candidates for extrema, not guarantees.
\end{remark*}

\begin{remark*}[Dependencies]
\begin{itemize}
  \item \hyperref[def:relative-extremum]{Definition: Relative Extremum}
  \item \hyperref[def:derivative-at-a-point]{Definition: Derivative at a Point}
  \item \hyperref[def:left-hand-derivative]{Definition: Left-Hand Derivative}
  \item \hyperref[def:right-hand-derivative]{Definition: Right-Hand Derivative}
  \item \hyperref[thm:differentiability-and-one-sided-derivatives]{Theorem: Differentiability and One-Sided Derivatives}
\end{itemize}
\end{remark*}

% ---------------------------------------------------------

\begin{theorem}[Necessary Condition for an Interior Extremum]
\label{thm:necessary-condition-extremum}
Let $I \subseteq \mathbb{R}$ be an open interval, let $f : I \to \mathbb{R}$, and let $c \in I$. If $f$ has a relative extremum at $c$, then either $f'(c)$ does not exist, or
\[
f'(c) = 0.
\]
\hyperref[prf:relative-extremum-derivatives-equal-0-or-dont-exist]{Go to proof.}
\end{theorem}

\begin{remark*}[Standard quantified statement]
\[
\operatorname{LocalExtremum}(f,c)
\;\Rightarrow\;
\lnot\operatorname{DifferentiableAt}(f,c) \;\lor\; \operatorname{DerivativeAt}(f,c,0).
\]
\end{remark*}

\begin{remark*}[Interpretation]
This is a weakening of the Interior Extremum Theorem that drops the differentiability hypothesis. Without assuming $f'(c)$ exists, the conclusion is a disjunction: either the derivative fails to exist at $c$, or it exists and equals zero. This is the correct form for identifying \emph{critical points} --- points that are candidates for extrema. The set of critical points is the union of points where $f'$ is zero and points where $f'$ does not exist.
\end{remark*}

\begin{remark*}[Dependencies]
\begin{itemize}
  \item \hyperref[thm:interior-extremum-theorem]{Theorem: Interior Extremum Theorem}
\end{itemize}
\end{remark*}

% ---------------------------------------------------------

\begin{corollary}[Necessary Condition for an Extremum]
\label{cor:necessary-condition-extremum}
Let $I \subseteq \mathbb{R}$ be an open interval, let $f : I \to \mathbb{R}$, and let $c \in I$. If $f$ has a relative extremum at $c$, then
\[
f'(c) = 0 \quad \text{or} \quad f'(c) \text{ does not exist.}
\]
\hyperref[prf:relative-extremum-derivatives-equal-0-or-dont-exist]{Go to proof.}
\end{corollary}

\begin{remark*}[Standard quantified statement]
\[
\operatorname{LocalExtremum}(f,c)
\;\Rightarrow\;
\operatorname{DerivativeAt}(f,c,0) \;\lor\; \lnot\operatorname{DifferentiableAt}(f,c).
\]
\end{remark*}

\begin{remark*}[Dependencies]
\begin{itemize}
  \item \hyperref[thm:necessary-condition-extremum]{Theorem: Necessary Condition for an Interior Extremum}
\end{itemize}
\end{remark*}
